\documentclass{like}

%\selectlanguage{english}
%\selectlanguage{german}



% BITTE mit pdf Latex oder Schnelles Übersetzen kompilieren sonst kommen Fehlermeldungen !!!!

\begin{document}



%%%%%%%%%%%%%%%%%%%%%%%%%%%%%%%%%%%%%%
%% Ab hier beginnt der Text
%%%%%%%%%%%%%%%%%%%%%%%%%%%%%%%%%%%%%%
  

%\normalsize
%Define our elements for our title page
\makeatletter

\renewcommand{\title}[1]{\gdef\@title{#1}}
\newcommand{\artderarbeit}[1]{\gdef\@artderarbeit{#1}}
\renewcommand{\author}[1]{\gdef\@author{#1}}
\newcommand\@matrikelnummer{\relax}
\newcommand{\matrikelnummer}[1]{\gdef\@matrikelnummer{#1}}
\newcommand{\professor}[1]{\gdef\@professor{#1}}
\newcommand\@studiengang{\relax}
\newcommand{\studiengang}[1]{\gdef\@studiengang{#1}}
\newcommand{\university}[1]{\gdef\@university{#1}}
\newcommand{\faulogo}[1]{\gdef\@faulogo{Abbildungen/#1}}
\newcommand{\likelogo}[1]{\gdef\@likelogo{Abbildungen/#1}}
\newcommand{\submitdate}[1]{\gdef\@submitdate{#1}}
\newcommand{\startdate}[1]{\gdef\@startdate{#1}}
\newcommand{\declaration}[1]{\gdef\@declaration{#1}}
\newcommand\@betreuerB{\relax}
\newcommand{\betreuerB}[1]{\gdef\@betreuerB{#1}}
\newcommand\@betreuerC{\relax}
\newcommand{\betreuerC}[1]{\gdef\@betreuerC{#1}}
\newcommand\@betreuerD{\relax}
\newcommand{\betreuerD}[1]{\gdef\@betreuerD{#1}}
 
%-----------------------------------------------------------------------
% Title page definition.
%-----------------------------------------------------------------------
\newcommand{\titlep}{
        \thispagestyle{empty}
        \begin{titlepage}
		\enlargethispage{4cm}
%        \null
		\vspace*{-2cm} 
        \begin{center}
{\rmfamily\large \textbf{ Friedrich-Alexander-Universität Erlangen-Nürnberg}}
\vskip 1,4cm
\begin{center}
%\includegraphics[scale=0.80, natwidth=1024, natheight=200]{\@faulogo} \hfill
\includegraphics[scale=0.80]{\@faulogo} \hfill
\vskip .8cm
{\large \bfseries Lehrstuhl für Informationstechnik\\(Schwerpunkt Kommunikationselektronik)}
\vskip 1.2cm
%\includegraphics[scale=0.40, natwidth=1024, natheight=200]{\@likelogo} \hfill
\includegraphics[scale=0.40]{\@likelogo} \hfill
\end{center}
\vskip .3cm
{\Large \@artderarbeit~mit dem Thema:\\}
\vskip 1.7cm
{\rmfamily\Large\bfseries\expandafter{\@title}}
        \end{center}
\vskip 1.5cm
		
\begin{table}[h!]
	\flushleft
%\setlength{\tabcolsep}{10pt}
	\large
\renewcommand{\arraystretch}{1,1}
	\begin{tabular}{p{4cm}l}
Bearbeiter     & {\rmfamily\large\expandafter{\@author}}\tabularnewline[.6ex]
\ifthenelse{\equal{\@matrikelnummer}{\relax}}{} {Matrikelnr. & \textrm{\@matrikelnummer}\tabularnewline[.6ex]}
\ifthenelse{\equal{\@studiengang}{\relax}}{}    {Studiengang & \textrm{\@studiengang}\tabularnewline[.6ex]}
Betreuer  & {\rmfamily\large\expandafter{\@professor}}\tabularnewline[.6ex]
                & \ifthenelse{\equal{\@betreuerB}{\relax}}{\vspace*{-0.6cm}}{\rmfamily\large{\@betreuerB}} \tabularnewline[.6ex]
                & \ifthenelse{\equal{\@betreuerC}{\relax}}{\vspace*{-0.9cm}}{\rmfamily\large{\@betreuerC}}  \tabularnewline[.6ex]
                & \ifthenelse{\equal{\@betreuerD}{\relax}}{\vspace*{-0.7cm}}{\rmfamily\large{\@betreuerD}} \tabularnewline[.6ex] Beginn         & {\rmfamily\large\expandafter{\@startdate}}\tabularnewline[.6ex]
Ende           & {\rmfamily\large\expandafter{\@submitdate}}            
	\end{tabular}
\end{table}
	        \vskip1cm
        \end{titlepage}
        \newpage}
        \makeatother
        
 


%-----------------------------------------------------------------------
% redefined \maketite
%-----------------------------------------------------------------------
\renewcommand{\maketitle}{%
                \titlep
	     }
%-----------------------------------------------------------------------
%  hier die eigenen Daten eintragen
%----------------------------------------------------------------------
    
\title{Latex für Abschlussarbeiten} % Mussfeld
\artderarbeit{Masterarbeit}         % Mussfeld
\author{Name Student}               % Mussfeld
\studiengang{Elektrotechnik-Elektronik-Informationstechnik}        % if Feld
\matrikelnummer{1234545646}         % if Feld
\faulogo{fau.png}                   % Mussfeld
\likelogo{like.png}                 % Mussfeld
\professor{Prof. Dr.-Ing. Jörn Thielecke}       % Mussfeld
\betreuerB{Markus Hiller, M.\,Sc.}           % if Feld
\betreuerC{Florian Particke, M.\,Sc.}              % if Feld
%\betreuerD{Betreuer D}             % if Feld  u.U. auskommentieren 
\startdate{01. Januar 2015}         % Mussfeld
\submitdate{08. August 2016}        % Mussfeld

% \maketitle

\maketitle
%BESTÄTIGUNG

\chapter*{Bestätigung}


\textbf{Erklärung:}\\

Ich versichere, dass ich die Arbeit ohne fremde Hilfe und ohne Benutzung anderer als der angegebenen Quellen angefertigt habe und, dass die Arbeit in gleicher oder ähnlicher Form noch keiner anderen Prüfungsbehörde vorgelegen hat und von dieser als Teil einer Prüfungsleistung angenommen wurde. Alle Ausführungen, die wörtlich oder sinngemäß übernommen wurden, sind als solche gekennzeichnet.
\vskip 2.5 cm
Erlangen, den (hier Datum eintragen)%{\rmfamily{\@submitdate}}
 \underline{\hspace{7cm} }
\vskip 1.5 cm

%DANKSAGUNG
%\pagestyle{scrheadings}
\chapter*{Danksagung}
%\headmark%chapter = Titel
%\section*{}
An dieser Stelle möchte ich mich bei all denjenigen bedanken, die mich während der Anfertigung dieser Masterarbeit unterstützt und motiviert haben.

Zuerst gebührt mein Dank Henrik Bey, der meine Masterarbeit betreut und begutachtet hat. Für die hilfreichen Anregungen und die konstruktive Kritik bei der Erstellung dieser Arbeit möchte ich mich herzlich bedanken.

Ebenfalls möchte ich mich bei meinen Kommilitonen Martin und Amin bedanken, die mir mit viel Geduld, Interesse und Hilfsbereitschaft zur Seite standen.

Abschließend möchte ich mich bei meinen Eltern Barbara und Matthias bedanken, die mir mein Studium durch ihre Unterstützung ermöglicht haben und stets ein offenes Ohr für meine Sorgen hatten.

Sebastian Weller,

Erlangen, 06.07.2018

 %Optional
%Thema und Aufgabenstellung
 
%\pagestyle{scrheadings}

\chapter*{Thema und Aufgabenstellung}
\markboth{Thema und Aufgabenstellung}{Thema und Aufgabenstellung}
\textbf{Thema:}\par\smallskip

 Modellfehler in optimierungsbasierter kombinierter Planung und Regelung für Rennwagen

\par\bigskip
\textbf{Aufgabenstellung:}\par\smallskip
\par\smallskip  
Die Automatisierung des Fahrens schließt sowohl die Planung als auch die Regelung des Fahrzeugs mit ein. Häufig werden beide Bestandteile hierarchisch voneinander getrennt. Dies ist sinnvoll, solange das kontrollierte Fahrzeug sicher innerhalb der Aktuatorlimitierungen betrieben werden soll, oder wenn die Trennung bereits durch die Problemstellung gegeben ist (Zieltrajektorie bereits vorgegeben) \cite{Williams2016AggressiveDW}.  

In anderen Fällen, z.B. wenn die gewünschte Dynamik wie in einer Rennsituation im Grenzbereich liegt, bietet sich eine kombinierte Planung und Regelung an. In diesem Beispiel würde die Kostenfunktion eine Minimierung der Rundenzeit beinhalten, während gleichzeitig die Beschränkungen des Fahrzeugs berücksichtigt werden.

Für derartige Probleme ist die modellprädiktive Regelung (MPC) bzw. eines ihrer Derivate besonders geeignet. Dabei kommt es immer zu einem sogenannten Modellfehler, der von der Komplexität und Genauigkeit des verwendeten Modells abhängt.

Das Ziel dieser Arbeit ist es, den Abfall bei der Leistung des Regelungsansatzes durch den Modellfehler zu untersuchen. Dafür soll eine Simulation verwendet werden.
Die Arbeit soll folgende Punkte beinhalten:
\begin{itemize}
\item Auswahl einer passenden Simulationsumgebung und deren Inbetriebnahme
\item Implementierung verschiedener (gegebener) Modelle für die Simulation
\item Implementierung des MPC-Ansatzes
\item Entwicklung einer einfachen Evaluationsmethode, um die Leistungsfähigkeit des Reglers zu untersuchen
\item Vergleich verschiedener Kombinationen aus Regler- und Simulationsmodellen
\end{itemize}
%Kurzzusammenfassung

\chapter*{Kurzzusammenfassung}

Mit zunehmender Rechenleistung und Erfahrung der Automobilbranche mit autonomen Fahrzeugen rückt auch das Thema der selbstfahrenden Rennautos immer mehr in den Fokus. Das ROBORACE Projekt ist hier Vorreiter mit seiner ausgefeilten Hardwareplattform und dem bereits in öffentlichen Events gezeigten Fahrleistungen. Auch die Formula Student (FS) verschließt sich nicht diesem Trend und hat 2017 die Rubrik Driverless ins Leben gerufen. \\
Diese Masterarbeit beschreibt einen Ansatz zur Echtzeitregelung und Trajektorienplanung für ein Rennauto das in diesem Wettbewerb antreten soll. Die Basis hierfür ist ein Model Predictive Control (MPC) Algorithmus. Er vereint die Regelung und Trajektorienplanung und ist adaptiv bezüglich verschiedener Fahrsituationen und Ziele.
Als Ausgangssituation wird angenommen, dass das Fahrzeug bereits eine Runde auf einem unbekannten Kurs absolviert und nun eine genaue Karte des Kurses errechnet hat. 
Um das MPC nutzen zu können, muss ein Fahrzeugmodell hinterlegt werden. Je genauer dieses ist, desto näher können die Grenzen des realen Fahrzeugs angenähert werden. Neben der Auslegung für das aktuellste FS-Fahrzeug des High Octane Motorsports Teams für die Driverless Umrüstung, wird untersucht, ab welchem Punkt ein kinematisches Modell nicht mehr ausreicht, um das Fahrzeug sicher auf dem Rennkurs zu halten.



%\abstract 
\chapter*{Abstract}
In recent years with more and more computational power and sophisticated algorithms, autonomous cars are close to becoming a reality. An ancillary effect of this development is the emerging of driverless racecars. For this idea to become a reality, three key ingredients have to be met. A very good approximation of the location of the racecar in a virtual map, a trajectory planning algorithm and a controller that will drive the car on the computed path. In this thesis the model predictive control(MPC) algorithm is used, to combine the trajectory planning and controlling for an autonomous race car, build for the formula student driverless competition. A simulation environment will be developed to evaluate the effectiveness of the approach. One key part of this thesis is to investigate how prone to model errors the MPC-approach is and what performance can be expected for the use of a kinetic bicycle model as system description in the (MPC)-algorithm.
%%Abkürzungsverzeichnis
%\pagestyle{scrheadings}
%\chapter*{Abkürzungsverzeichnis}

\chapter*{Abkürzungsverzeichnis}
\addcontentsline{toc}{chapter}{Abkürzungsverzeichnis} %Abkürzungsverzeichnis kommt ins InhaV.





\begin{acronym}[SEPSEP] 
 \acro{Abk.}{Abkürzung}
 \acro{z.B.}{zum Beispiel}
 \acro{MPC}{Model Predictive Control}
 \acro{FS}{Formula Student}
 \acro{FSG}{Formula Student Germany}
 \acro{JIT}{just-in-time}
 \acro{SLSQP}{Sequential Least SQuarez Programming}
 \acro{SFML}{Simple and Fast Multimedia Library}
\end{acronym}


 (nach hinten verschoben)

\tableofcontents % erstelle hier das Inhaltsverzeichnis
\markboth{Contents}{Contents}
%\listoffigures % erstelle hier das Abbildungsverzeichnis 
%\listoftables % erstelle hier das Tabellenverzeichnis
\enlargethispage{2cm} %Bitte 2x kompilieren - sonst erscheint es nicht


\pagenumbering{arabic} %Seitennummerierung arabisch
\setcounter{page}{1}   %1. Kapitel soll mit Seitenzahl 1 beginnen
%------------------------------------------------------------------
\chapter{Einleitung}
\section{Gliederung}
Befehle zum Gliedern von Texten:\par
Nummeriert und im InhV.:\nopagebreak\par\medskip\nopagebreak

\begin{minipage}{0.4\textwidth}
\begin{verbatim}
\chapter{Kapitel - Ebene 1 (1 Einleitung)}
\section{Abschnitt - Ebene 2 (1.1 Gliederung}
\subsection{Unterabschnitt - Ebene3 1.1.1....} 
\subsubsection{Unterunterabschnitt - Ebene 4 (nicht sinnvoll)}
\end{verbatim}
\end{minipage}


\subsection*{nicht nummerierte Überschrift}
Alles, was nicht ins Inhaltsverzeichnis soll ist mit einem *
 versehen: -> steht für Kapitel, die nicht nummeriert sind und somit auch nicht im Inhvz. stehen\par
Befehle: \par

\begin{minipage}{0.2\textwidth}
\begin{verbatim}
\chapter*{Thema und Aufgabenstellung}
\section* {Sektion/Abschnitt} unnummiert 
\subsection*{Unterabschnitt - Ebene3 } 
\subsubsection*{Unterunterabschnitt - Ebene 4}
 \end{verbatim}
 \end{minipage}


\section{Text formatieren}

Folgende Befehlen erzeugen einen Zeilenumbruch:\\ (1ex entspricht der Höhe des Buchstaben \glqq x\grqq\ in der aktuellen Schriftart)\par\medskip
\qquad\begin{minipage}{0.8\textwidth}
\begin{verbatim}
\\
\newline
\\[1ex]
\end{verbatim}
\end{minipage}\par\bigskip

Folgende Befehle erzeugen einen neuen Absatz:

\qquad\begin{minipage}{0.8\textwidth}
\begin{verbatim}
\par
\par\smallskip
\par\medskip
\par\bigskip
\end{verbatim}
\end{minipage}\par\bigskip

Folgende Befehle erzeugen einen zusätzlichen horizontalen Abstand:\nopagebreak\par\medskip\nopagebreak

\qquad\begin{minipage}{0.8\textwidth}
\begin{verbatim}
\,                 % kleiner Abstand
\;                 % mittlerer Abstand
\quad              % großer Abstand
\qquad             % sehr großer Abstand
\!                 % kleiner negativer Abstand
\end{verbatim}
\end{minipage}\par\bigskip


Folgende Befehle zum Seitenumbruch:\nopagebreak\par\medskip\nopagebreak



\section{Umbrüche}
Der Text wird normalerweise automatisc umbrochen. Zum expliziten Beenden einer Zeile benutzen Sie
\begin{verbatim}
\\[abstand] % bzw.
\\[2pt]
\end{verbatim}
Beispiel\\[4pt]
Wenn mann zum Beispiel eine Überschrift
\begin{verbatim}
\\in zwei Zeilen umbrechen möchten,
\end{verbatim}
bitte ein
\begin{verbatim}
\protect
\end{verbatim}
benutzen. Zusätzlich kann an einer so beendeten Zeile noch gleichzeitig ein Seitenumbruch verhindert werden,
indem man
\begin{verbatim}
\\*[abstand]
\end{verbatim}
eingibt.
Wenn man Latex einen Spielraum lassen will: Dafür dient der Befehl
\begin{verbatim}
\linebreak[zwang]
\end{verbatim}

 Für die Dringlichkeit des-Befehls können Sie die Zahlen von 0 bis 4 eintragen,wobei 0 eine milde Wunschäußerung darstellt und 4 das Zeilenende erzwingt. Gleichzeitig werden beim Beenden einer Zeile die Wörter gestreckt und dabei gleichmäßig überdie Zeile verteilt.


\begin{verbatim}
Der Zwerg nahm seine Axt und hieb mehrere Äste ab.\\
Der Zwerg nahm seine Axt und hieb mehrere Äste ab.\linebreak[4]
Der Zwerg nahm seine Axt und hieb mehrere Äste ab.\linebreak[2]
Der Zwerg nahm seine Axt und hieb mehrere Äste ab.\linebreak[3]
Sie banden sie mit Bogensehnen zusammen
\end{verbatim}
Der Zwerg nahm seine Axt und hieb mehrere Äste ab.\\
Der Zwerg nahm seine Axt und hieb mehrere Äste ab.\linebreak[4]
Der Zwerg nahm seine Axt und hieb mehrere Äste ab.\linebreak[2]
Der Zwerg nahm seine Axt und hieb mehrere Äste ab.\linebreak[3]
Sie banden sie mit Bogensehnen zusammen
\par
oder:
\begin{verbatim}
\nolinebreak[zwang]
\end{verbatim}
\section{Schriftart formatieren}

Folgende Befehle für die Schriftgröße:\par\smallskip

\qquad\begin{minipage}{0.4\textwidth}
\vspace*{-0,3cm}
\begin{verbatim}
\tiny
\scriptsize
\footnotesize
\small
\normalsize
\large
\Large
\LARGE
\huge
\Huge
\end{verbatim}
\end{minipage}
\hspace*{0,8cm}
\begin{minipage}{0.4\textwidth}
\tiny tiny\,
\scriptsize scriptsize\,
\footnotesize footnotesize\;
\small small\;
\normalsize normalsize\\
\large large\;
\Large Large\
\LARGE LARGE\\
\huge huge\,
\Huge Huge\;
\end{minipage}\normalsize


Folgende Befehle für die Schriftart :

\qquad\begin{minipage}{0.2\textwidth}
\begin{verbatim}
\textrm{} 
\texttt{} 
\textbf{} 
\textit{} 
\textsc{} 
\textsl{} 
\end{verbatim}
\end{minipage}\quad
\begin{minipage}{0.4\textwidth}
\textrm{Textbeispiel} \\
\texttt{Textbeispiel} \\
\textbf{Textbeispiel} \\
\textit{Textbeispiel} \\
\textsc{Textbeispiel} \\
\textsl{Textbeispiel} 
\end{minipage}

\par\bigskip


Folgende Befehle zum Hervorheben (Argument):%\nopagebreak\par\medskip\nopagebreak

\begin{minipage}{0.2\textwidth}
\begin{verbatim}
\emph{}
\underline{}
\end{verbatim}
\end{minipage}\quad
\begin{minipage}{0.4\textwidth}
\emph{Textbeispiel}\\
\underline{Textbeispiel}
\end{minipage}
\subsection{Default Fonts}
Für das gesamte Dokument wurde eine Schriftart im cls.-File generiert\par\nopagebreak
\begin{minipage}[ht!]{0.4\textwidth}
mittes des Befehls->\\
und dem Paket Times->\\
\end{minipage}
\begin{minipage}[ht!]{0.4\textwidth}
\vspace*{-.6cm}
\begin{verbatim}
\addtokomafont{disposition}{\rmfamily}
\RequirePackage{times} 
\end{verbatim}
\end{minipage}


\chapter{Grundlagen}


\section{Tabellen}

Folgende Befehle für eine Tabelle:%\nopagebreak\par\medskip\nopagebreak


\qquad\begin{minipage}{0.8\textwidth}
\begin{verbatim}
&                 % neue Spalte
\\                % neue Zeile
\\[1ex]           % neue Zeile mit Abstand 1ex
\hline            % horizontale Linie
\end{verbatim}
\end{minipage}\par\bigskip

Ein Beispiel:
Syntax:
\begin{verbatim}
\begin{tabular}[Position]{Spalten} ... \end{tabular}
\end{verbatim}


Parameter:

Position
    (optional) b oder t. 
Spalten
    eine Spaltendefinition.\\ 
Beispiel für eine Tabelle:
 

\begin{minipage}[t!]{0.5\textwidth}
  \begin{tabular}[t]{lrr@{.}l}
    Position 1 & DM &  5 & 49\\
    Position 2 & DM & 10 & 95\\
   \end{tabular}

\end{minipage}\hspace*{5pt}
\begin{minipage}[t!]{0.5\textwidth}

\begin{verbatim}

  \begin{tabular}[t]{lrr@{.}l}
    Position 1 & DM &  5 & 49\\
    Position 2 & DM & 10 & 95\\
  \end{tabular}
  \end{verbatim}
\end{minipage}
\section{Einbindung von Tabellen}
Tabellen und Grafiken werden nach Kapitel nummeriert
 \begin{table}[h]
 	 \caption{Tabelle}
  	\begin{tabular}{rl}
    linke Spalte & rechte Spalte
  	\end{tabular}
 \end{table}
 
 \begin{verbatim}
 \begin{table}[h]
 	 \caption{Tabelle}
  	\begin{tabular}{rl}
    linke Spalte & rechte Spalte
  	\end{tabular}
 \end{table}
 \end{verbatim}
Mehr dazu : \\
\url{http://www.weinelt.de/latex/tabular.html}\\
\url{http://www.weinelt.de/latex/index.html}\\
\url{http://www.torsten-schuetze.de/tex/tabsatz-2004.pdf}
\section{Figures}
Figures sind Umgebungen mit deren Hilfe man beliebige Objekte wie zum Beispiel
Tabellen oder Graphiken im laufenden Text platzieren kann.\nopagebreak\par\medskip\nopagebreak
\qquad\begin{minipage}{0.8\textwidth}
\begin{verbatim}
\begin{figure}[ht!]
  \caption{Bildüberschrift}
  % beliebiges Objekt
\end{figure}
\end{verbatim}
\end{minipage}\par\bigskip
Die in eckigen Klammern angegebenen Parameter geben die Positionierung der Figure-Umgebung an \nopagebreak\par\medskip\nopagebreak
\qquad\begin{minipage}{0.8\textwidth}
\begin{verbatim}
h       % Stelle, an der sich Figure befindet
t       % Anfang der (unter Umständen folgenden) Seite
b       % Ende der Seite
p       % Auf einer extra Seite
!       % Parameter haben Priorität vor Latex-internen
          Parametern
\end{verbatim}
\end{minipage}\par\bigskip
Mit dem Befehl \verb+\caption{}+ vor dem Bild erzeugt man eine Figure-Übeschrift, bzw nach dem Objekt eine Figure-Unterschrift

\section{Graphiken}
Um Grafiken einbinden zu können ist das Paket \verb+graphicx+ in der Präambel eingebunden:
\qquad\begin{minipage}{0.8\textwidth}
\begin{verbatim}
\RequirePackage[dvips]{graphicx}
\end{verbatim}
\end{minipage}\par\bigskip
Standardmäßig sollte der Parameter \verb+dvips+ verwenden werden, für \glqq PDFLatex\grqq \, ist \verb+pdftex+ empfehlenswert.
Der eigentliche Befehl zum Einbinden von Grafiken lautet:
\subsection{Beispiel für Grafiken}

\qquad\begin{minipage}{0.8\textwidth}
\begin{verbatim}
\includegraphics[key=values, key=values]{file}
\end{verbatim}
\end{minipage}\par\bigskip
\verb+key+ kann dabei die folgenden Parameter annehmen die, mit einem Komma getrennt, beliebig kombiniert werden können:\nopagebreak\par\medskip\nopagebreak
\qquad\begin{minipage}{0.8\textwidth}
\begin{verbatim}
width             % Breite festlegen
height            % Höhe festlegen
angle             % Rotation gegen den Uhrzeigersinn
scale             % Skalierung 
\end{verbatim}
\end{minipage}\par\bigskip

 \subsection{Einbindung von 2 Grafiken nebeneinander}
  \begin{figure}[ht!]
  \begin{centering}
  \subfigure[Bezeichnung der linken Grafik]{\includegraphics[width=0.33\textwidth]{Abbildungen/mehrkeulenantenne_prinzip.png}}
  \hspace*{1cm}
    \subfigure[Bezeichnung der rechten Grafik]{\includegraphics[width=0.33\textwidth]{Abbildungen/pcb_ant.jpg}}
 \caption{Bildunterschrift}
  \end{centering}
  \end{figure}
\subsection{Darstellung von 3 Bildern nebeneinander}
mit einer Minipage Umgebung

 \begin{figure}[h!]
\begin{minipage}[t]{0.3\textwidth}
\includegraphics[width=\textwidth]{Abbildungen/pcb_ant.jpg}
Bild 1
\end{minipage}
\hspace*{0.2cm}
\begin{minipage}[t]{0.3\textwidth}
\includegraphics[width=\textwidth]{Abbildungen/motion-sensor-pcb.jpg}
Bild 2
\end{minipage}
\hspace*{0.2cm}
\begin{minipage}[t]{0.3\textwidth}
\includegraphics[width=\textwidth]{Abbildungen/peilantenne_c-band_bottom.png}
Bild 3
\end{minipage}
\caption{Bildunterschrift Unterschrift Bild x Die auf die Rotationsfrequenz des Innenzylinders normierten Eigenfrequenzen der gefundenen Grundmoden der Taylor-Strömung für h}
  \end{figure}




\section{Bildunterschriften}
 \subsection*{Einbindung Grafik}
  \begin{figure}[h!]
  \begin{centering}
  {\includegraphics[width=0.33\textwidth]{Abbildungen/fau.png}}
   \caption{Unterschrift Bild x Die auf die Rotationsfrequenz des Innenzylinders normierten Eigenfrequenzen der gefundenen Grundmoden der Taylor-Strömung für h(Die azimutale Wellenzahl ist mit m  bezeichnet.)}
  \end{centering}
  \end{figure}
%Bildunterschriften
Mit dem Paket
\begin{verbatim}
\usepackage{ragged2e}
\usepackage{caption}
\end{verbatim}
sowie den Einstellungen:
\begin{verbatim}
\captionsetup{format=plain,indention=2.3cm,font=small,labelfont=bf,
labelsep=colon,justification=RaggedRight} ->linksbündig mit 
indention=Einzug
Ansonsten: wie hier:
\captionsetup{font=small,labelfont=bf}
-> Blocksatz
\end{verbatim}

Sollte es wider erwarten zu komischen Trennungen kommen - kann man die Befehle 
\begin{verbatim}
\sloppy % weniger Worttrennungen, größere Wortabstände
\fussy % viele Worttrennungen, "schönere" Wortabstände
\end{verbatim}
einfügen 

\section{Aufzählungen}
\par

\begin{minipage}[t!]{0.38\paperwidth}
\begin{verbatim}
\begin{itemize}
	\item  test 2. ebene
	\item  test 2. ebene
		\begin{itemize}
			\item  test 3. ebene
			\item  test 3. ebene
				\begin{itemize}
					\item  test 4. ebene
					\item  test 4. ebene
				\end{itemize}
		\end{itemize}
	\end{itemize}
\end{itemize}
Nummerierte Listen:
	\begin{enumerate}
	\item Textbeispiel
		\begin{enumerate}
		\item Textbeispiel
		\item Textbeispiel
			\begin{enumerate}
		\item Textbeispiel
		\item Textbeispiel
\end{enumerate}
\end{enumerate}
\end{verbatim}

\end{minipage}\quad
\begin{minipage}[t!]{0.38\textwidth}
\vspace*{.5cm}
\begin{itemize}
\item 1. Ebene
\begin{itemize}
\item  test 2. ebene
\item  test 2. ebene
\begin{itemize}
\item  test 3. ebene
\item  test 3. ebene
\begin{itemize}
\item  test 4. ebene
\item  test 4. ebene
\end{itemize}    
\end{itemize}
\end{itemize}
\end{itemize}(alle Ebenen wieder schließen!)
\par\bigskip
\begin{enumerate}
	\item Ebene 1
		\begin{enumerate}
			\item Ebene 2
			\item Ebene 2
				\begin{enumerate}
					\item Ebene 3
					\item Ebene 3
				\end{enumerate}
		\end{enumerate}
		\item Ebene 1
	\end{enumerate}
\end{minipage}
\par






\section{Farben}
Um Farben verwenden zu können muss das Paket \verb+color+ Header eingebunden werden. 
Folgende Befehle dienen dem Ändern von Farben: \nopagebreak\par\medskip\nopagebreak
\begin{minipage}{0.8\textwidth}
\begin{verbatim}
\color{color}          % Ändert die Textfarbe ab hier
\color{color}{Text}    % Ändert die Textfarbe von "Text".
\pagecolor{color}      % Ändert die Hintergrundfarbe
\colorbox{color}{Text} % Unterlegt "Text" mit der gewünschten Farbe
\fcolorbox{color1}{color2}{Text} 
% Unterlegt "Text" mit color1 und Rand color2 \end{verbatim}
 Folgende Farben sind definiert:
 \begin{verbatim}
\definecolor{color0}{rgb}{0,0,0}% black
\definecolor{color1}{rgb}{0.22,0.45,0.70}% light blue
\definecolor{color2}{rgb}{0.45,0.45,0.45}% dark grey

\end{verbatim}
\end{minipage}\par\bigskip




\section{Matheumgebung}

Folgende Befehle, um die Matheumgebung einzuschalten:\nopagebreak\par\medskip\nopagebreak

\begin{minipage}{0.8\textwidth}
\begin{verbatim}
$             % Anfang/Ende der Matheumgebung im laufenden Text

\(            % Anfang der Matheumgebung im laufenden Text
\)            % Ende der Matheumgebung im laufenden Text

Einzeilig:
$$            % Anfang/Ende der Matheumgebung abgesetzter Formel 
				

\[            % Anfang der Matheumgebung in abgesetzter Formel 
\]            % Ende der Matheumgebung in abgesetzter Formel 

\begin{equation}  % Anfang der Matheumgeb. (abges.Formel)
\end{equation}    % Ende der Matheumgebung in abgesetzter Formel 

Mehrzeilig:

\begin{eqnarray}  % Anfang der Matheumgebung in abgesetzter Formel
\end{eqnarray}    % Ende der Matheumgebung in abgesetzter Formel
\end{verbatim}
\end{minipage}\par\bigskip


	\section{Formelnummerierung nach section}
	
	\begin{verbatim}
	wer dies andern möchte : 
	im cls-file
	\numberwithin{equation}{section} <- section in chapter umändern
	\end{verbatim}
	
		\begin{equation} 
			x^2 - 5 x + 6 = 0
		\end{equation}
		
		\begin{equation}
			x_1 = \frac{5 + \sqrt{25 - 4 \times 6}}{2} = 3
		\end{equation}

		\section{Nummerierung der Formeln}
  		 \begin{equation}
    	  E=mc^2
   		\end{equation}
\chapter{Implementierung von xy} 

	\section{Running Header}
	Wenn keine
	\begin{verbatim}
 \section 
 \end{verbatim}
angegeben ist, wird als Running Header nur 
das\begin{verbatim}
\chapter 
\end{verbatim} ausgegeben
Code dafür im cls.-file unter
\begin{verbatim}
\RequirePackage[headsepline]{scrpage2}
\pagestyle{scrheadings} 
\automark[chapter]{chapter}
\clearscrheadfoot 

Fonts:
 \setkomafont{pageheadfoot}{%Kopfzeilen in Kapitälchen
\normalsize\scshape}        
Fußzeile: mittig
\cfoot[\pagemark]{\pagemark}                      
\end{verbatim}
		%\section{Ich bin der auch der Running Header}
		
	
\section{ich bin der Running Header}
	




\chapter{Konzept}


	\section{Zitate für die Bibliography}
	\begin{verbatim}
Citation of  Adler ~\cite{Adler2016}.
Citation of Adler b\cite{Adler2016b}
	\end{verbatim}
	\subsection{Beispiel für Zitate}
% Zitate :
Citation of  Adler ~\cite{Adler2016}.
Citation of Adler b\cite{Adler2016b}

\chapter{Ergebnisse}

communicated with simplicity and charm. And in such a way that the reader will read on. (After all, that’s a reader’s job: to read, isn’t it?) And by the time your readers have reached this point in the finished copy, you will have convinced them that you not only respect their intelligence, but you also unders




\chapter{Diskussion}
Text beginnt hier...his, of course, is not the real copy for this advertisement. The real words will be written once you have approved the headline. Rest assured, the words will expand the concept. With clarity. Conviction. And even a little wit. Because in today’s competitive marketing environment, the body copy of your advertisment must lead the reader through a series of disarmingly 
\section{Test für Abstand der Sektion} 
simple thoughts. All your \ac{sup.}arguments must be communicated with simplicity and charm. And in such a way that the reader will read on. (After all, that’s a reader’s job: to read, isn’t it?) And by the time your readers have reached this point in the finished copy, you will have convinced them that you not only respect their intelligence, but you also understand their needs as consumers. As a result of which, your advertis
\subsection{Test für den Subsectionabstand}
simple thoughts. All your supporting arguments must be communicated with simplicity and charm. And in such a way that the reader will read on. (After all, that’s a reader’s job: to read, isn’t it?) And by the time your readers have reached this point in the finished copy, you will have convinced them that you not only respect their intelligence, but you also understand their needs as consumers. As a result of which, your advertis

\chapter[Zusammenfassung]{Zusammenfassung und Ausblick \\\hspace{5mm}( für mehrzeilige Titel)}
Texte, die eine \ac{Abk.} enthalten, die nicht im Duden steht, sollten ein \acl{Abk.}sverzeichnis haben. Eine \ac{Abk.}, die nicht verwendet wird, steht auch nicht im \acl{Abk.}sverzeichnis.



\appendix

%\chapter*{Anhang}
%\large{
%Anhang 1\par
%Anhang 2


% setzt die Bildunterschriften auf A1 (A=Anhang)
\renewcommand{\thefigure}{A\arabic{figure}} 
\chapter{Anhang}

\subsection*{Elektronischer Anhang }

\subsection*{Fahrzeugdaten}
\begin{figure}[hb!]
	\caption{Engine Power}
	\includegraphics[width=350pt]{Abbildungen/Engine_power.png}
	\label{fig:enginePower}
\end{figure}

\begin{figure}[hb!]
	\caption{Max Tire Force}
	\includegraphics[width=350pt]{Abbildungen/Engine_power.png}
	\label{fig:maxTireForce}
\end{figure}



 \subsection*{Einbindung Grafik im Anhang}
  \begin{figure}[h!]
  \begin{centering}
  {\includegraphics[width=0.33\textwidth]{Abbildungen/fau.png}}
   \caption{Unterschrift Bild x Die auf die Rotationsfrequenz des Innenzylinders normierten Eigenfrequenzen der gefun-denen Grundmoden der Taylor-Strömung für h(Die azimutale Wellenzahl ist mit m  bezeichnet.)}
  \end{centering}
  \end{figure}

\subsection{Code}

\begin{lstlisting}
[caption={ein paar Zeilen code}\label{lst:vehicleModel},captionpos=t] 
    for i in 0: N-1
    
    @NLexpression(m, x_dd, (VehicleModel.max_long_acc * x[8*i + 7]/10.0))
    
    @NLconstraints(m, begin
    x[(i + 1)*8 + 1] - (x[i * 8 + 1] + x[8*i + 3]*
    dt*cos(x[i*8 + 4] + atan(lr/(lf + lr) * tan(x[i*8 + 8])))) == 0
    x[(i + 1)*8 + 2] - (x[i * 8 + 2] + x[8*i + 3]*
    dt*sin(x[i*8 + 4] + atan(lr/(lf + lr) * tan(x[i*8 + 8])))) == 0
    x[(i + 1)*8 + 3] - (x[i * 8 + 3] + x_dd*dt) == 0
    x[(i + 1)*8 + 4] - (x[i * 8 + 4] + x[8*i + 3]*
    dt / lr*sin(atan(lr/(lf + lr) * tan(x[i*8 + 8])))) == 0
    atan(0.5 * (lf + lr) * mpc_struct.beta_max / x[i*8 + 3]^2)
    - atan(lr/(lf + lf) * tan(x[i*8 + 8])) >= 0  #max_beta - beta
    atan(0.5 * (lf + lr) * mpc_struct.beta_max / x[i*8 + 3]^2)
    + atan(lr/(lf + lf) * tan(x[i*8 + 8])) >= 0  #max_beta + beta
    end)
    end
\end{lstlisting}







%Abkürzungsverzeichnis
%\pagestyle{scrheadings}
%\chapter*{Abkürzungsverzeichnis}

\chapter*{Abkürzungsverzeichnis}
\addcontentsline{toc}{chapter}{Abkürzungsverzeichnis} %Abkürzungsverzeichnis kommt ins InhaV.





\begin{acronym}[SEPSEP] 
 \acro{Abk.}{Abkürzung}
 \acro{z.B.}{zum Beispiel}
 \acro{MPC}{Model Predictive Control}
 \acro{FS}{Formula Student}
 \acro{FSG}{Formula Student Germany}
 \acro{JIT}{just-in-time}
 \acro{SLSQP}{Sequential Least SQuarez Programming}
 \acro{SFML}{Simple and Fast Multimedia Library}
\end{acronym}




%\bibliographystyle{IEEEtran}
\bibliographystyle{alpha}
\bibliography{Inhalt/literatur}

% Bitte noch Index anhängen (kommt unter Linux zu Fehlermeldungen)
%\input{Inhalt/index}

%bitte Lebenslauf anhängen !!!
% kommt unter Linux zu Fehlermeldungen

%%%%%%%%%%%%%%%%%%%%%%%%%%%%%%%%%%%%%%%%%%%%%%%%%%%%%
          %   \begin  CV und Beginn des Dokuments
%%%%%%%%%%%%%%%%%%%%%%%%%%%%%%%%%%%%%%%%%%%%%%%%%%%%
%\enlargethispage{3cm}

\chapter*{Albert Einstein}
\markboth{Albert Einstein}{Albert Einstein} % falls der CV über 2 Seiten geht um einen "Running Header" zu haben


\flushright
% Code mit Bild:
%\includegraphics[scale=.2]{Abbildungen/albert-einstein.jpg}
%\vspace*{-5cm}{
%Code ohne Bild:
%\vspace*{-1cm}{
\subsection*{Persönliche Daten}
\flushleft
%\tabular-Umgebung
\normalsize
\begin{tabular}{lcl}
Adresse             & ~ & An der Kühruh 13\\
                    & ~ & 96123 Litzendorf\\[5pt]
Mobil               & ~ & 0170 - 9732890\\
Email               & ~ & sebastian.weller01@gmail.com\\[5pt]
Geburtsdatum        & ~ & 01.04.1992\\
Staatsangehörigkeit & ~ & deutsch\\
\end{tabular} 

\subsection*{Bachelorarbeit (optionaler Punkt)}
\begin{tabular}{lcl}
01/2016-07/2016     & ~~~~~ & Bachelorarbeit \\
                    & ~~~~~ & Ich bin das Thema der Bachelorarbeit\\[5pt]

\end{tabular} 
\\*  % Zeilenumbruch OHNE Seitenwechsel

\nopagebreak
\subsection*{Studium und Schulbildung}
\begin{tabular}{lcl}
01/2016 - 07/2016     & ~~~ &  Friedrich-Alexander-Universität Erlangen-Nürnberg\\
                      & ~~~ &  Hauptfächer Prokrastination und Bummelei\\
01/2010 - 01/2016     & ~~~ &  Albert-Einstein-Gymnasium, Erlangen\\
                      & ~~~ &  Leistungskurse: Feiern und Relaxen\\
\end{tabular}

\subsection*{Berufliche Erfahrungen / Praktika}
\begin{tabular}{lcl}
01/2016 - 07/2016     & ~~~ &  Wissenschaftlicher Hilfsmitarbeiter am Fraunhofer IIS\\
01/2016 - 07/2016     & ~~~ &  Praktikum bei Siemens Erlangen\\
\end{tabular}
\subsection*{Zusatzqualifikationen - (optional)}
\begin{tabular}{lcl}
Sprachen            &  & Deutsch (Muttersprache)\\
                    &  & Englisch (fließend in Wort und Schrift)\\
                    &  & Französisch (Grundkenntnisse)\\[5pt]
Programmiersprachen &  & Java \\
\end{tabular} 
\bigskip
\vspace*{.5cm}




Erlangen, den (Datum eintragen)\\*[20pt]\nopagebreak
\vspace{0.5cm}


\rule{5cm}{0.4pt}\\*\nopagebreak


Albert Einstein\\*\nopagebreak







\end{document}



 
